% !PRINTED 157
\weberpage{157}

Dieser Ausdruck zeigt nun, dass $R$, abgesehen von dem
Factor $(-1)^{mn}$, in der gleichen Weise von $\varphi(x)$ abhängt, wie
von $f(x)$, und dass wir also auch setzen können

\weq{7}{(-1)^{mn}\, R = b_0^n\, f(\beta_1)\, f(\beta_2)\, \ldots\, f(\beta_m).}

Die Discriminanten sind hiernach ein specieller Fall der
Resultanten, den man erhält, wenn man für $\varphi(x)$ die Derivirte
$f'(x)$ nimmt.

% [corrected 2026-08-07, 10x zoom strips 3/10 and 4/10] The whole theorem statement is
% letterspaced on the page, across the display and through "eine gemeinsame Wurzel haben." --
% the contrast against "f'(x) nimmt." on the line above is unmistakable. The transcription pass
% captured none of it. \gesperrt cannot span a display, so the run is wrapped in two pieces.
% Raised at high confidence by this page's translation pass and confirmed.
\gesperrt{Das Verschwinden der Resultante von $f(x)$ und
$\varphi(x)$ ist die nothwendige und hinreichende Bedingung
dafür, dass die beiden Gleichungen}

\[ f(x) = 0 \quad \text{und} \quad \varphi(x) = 0 \]

\gesperrt{eine gemeinsame Wurzel haben.}

% [corrected 2026-08-07, same strip] "Endgleichung" (both occurrences) and "Elimination" are
% letterspaced here -- Weber's habit of setting a term in Sperrung where he defines it.
Die Gleichung, die durch Nullsetzen der Resultante entsteht,
wird auch die \gesperrt{Endgleichung} aus $f(x) = 0$ und $\varphi(x) = 0$
genannt und die Aufstellung der \gesperrt{Endgleichung} die \gesperrt{Elimination}
von $x$.

Zur Berechnung der Resultanten kann man den Algorithmus
des grössten gemeinsamen Theilers anwenden (\S.\,6). Ist nämlich
$m \leqq n$, so kann man eine ganze Function $Q$ und eine ebensolche
Function $\psi(x)$, deren Grad kleiner ist als $m$, so bestimmen,
dass

\[ f(x) = Q\varphi(x) + \psi(x), \]

also, wenn man für $x$ eine der Wurzeln $\beta_i$ setzt,

\[ f(\beta_i) = \psi(\beta_i), \]

also

\weq{8}{b_0^n \overset{i}{\Pi} f(\beta_i) = b_0^n \overset{i}{\Pi} \psi(\beta_i);}

ist $m'$ der Grad von $\psi(x)$, so ist

\weq{9}{b_0^{m'} \overset{i}{\Pi} \psi(\beta_i) = R'}

die Resultante von $\varphi(x)$ und $\psi(x)$ und daher

\weq{10}{R = b_0^{n-m'}\, R'.}

Die Bildung der Resultante $R$ ist hierdurch auf die Bildung
von $R'$ zurückgeführt. Man kann nun $\varphi(x)$ wieder durch $\psi(x)$
theilen, und führt so die Bildung von $R'$ auf die Bildung einer
Resultante von Functionen immer niedrigeren Grades zurück, bis
man schliesslich auf eine Constante, d.\,h. eine Function der
Coëfficienten $a$, $b$ kommt, die mit der gesuchten Resultante identisch
sein muss.
